\chapter{电场中的导体与电介质}
\setcounter{example_counter}{1}

\section{电场中的导体}

\subsection{导体的静电平衡条件}

\subsubsection{静电平衡}

导体内部具有自由电子，将导体放在电场中，自由电子将在电场的作用下移动。

导体自由电子的移动，将使导体出现\textbf{感生电荷}，这些电荷也产生电场，与外电场抵消。

当导体内部和表面都没有电荷定向移动时，称为\textbf{静电平衡}。

\subsubsection{静电平衡条件}

处于静电平衡的导体，内部电场为零，表面紧邻处的电场强度与导体表面垂直。

处于静电平衡的导体是等势体，其表面是等势面。

\note{此处所说的电场，应理解为外电场和导体电荷移动所产生电场的叠加。}

\subsubsection{电荷分布}

\begin{itemize}
    \item 处于静电平衡的导体，内部各处净电荷为零，电荷只能分布在表面
    \item 表面各处的面电荷密度，与表面紧邻处的电场强度（合场强）大小成正比
    \item 对于孤立的导体，面电荷密度与表面曲率有关，曲率越大则面电荷密度也越大
\end{itemize}

\note{尖端附近的面电荷密度最大，可产生\textbf{尖端放电}现象。}

利用高斯定理，可求出导体表面附近的电场强度

\begin{formula}
    $E=\dfrac{\sigma}{\varepsilon_0}$
\end{formula}

\warn{注意无限大平板的电场强度大小为$E=\dfrac{\sigma}{2\varepsilon_0}$，不要混淆。}

\subsubsection{静电屏蔽}

金属壳外的电荷在壳内产生的电场为零，即金属壳屏蔽壳外电荷对壳内的影响。

\warn{注意金属壳不能屏蔽壳内电荷对壳外的影响。}

\begin{example}{静电感应}
    (1)~带电量为$+Q$的金属球在P点产生的电场强度大小为$E$，现在P点放一个正试探电荷，它的受力大小等于$qE$吗？如果是负试探电荷呢？(2)~将无限大金属板放在匀强电场$\vec{E}$中，求金属板表面的面电荷密度。
    
    【补充练习】半径为$R$的导体球外部有一电荷$Q$，它到球心的距离为$r>R$。(1)~求球心的场强和电势；(2)~若将导体球接地，求导体球的带电量。
\end{example}

\subsection{有导体存在的电场分析}

\subsubsection{无限大平板}

对于无限大金属平板，无论多薄，都不能忽略其厚度，需要分别设出两个表面的面电荷密度$\sigma$

然后利用每个面的场强为$E=\dfrac{\sigma}{2\varepsilon_0}$和平板内部电场强度为零，利用高斯定理进行求解。

若有某个面接地，则该面上没有电荷，该平板的电势为零。

\begin{example}{无限大导体平板}
    A、B为两个面积$S$足够大的导体平板，平行放置，间距为$d$，分别带电荷$+Q_1,+Q_2$。(1)~求A、B之间的电场强度和电势差；(2)~若将B的下表面接地，求A、B之间的电场强度和电势差。

    【补充练习】三个平行金属板A、B、C的面积都是$S$（足够大），间距分别为$d_1,d_2$，A、C两板接地，B板带电$+Q$。(1)~求各金属板表面的电荷面密度；(2)~求B板的电势。
\end{example}

\subsubsection{导体球与球壳}

若在导体球或球壳外部有电荷，则感应电荷只出现在外表面，球或球壳内部电场强度为零。

若在导体球壳内部有电荷，则感应电荷会出现在内外表面，内表面电荷量与内部电荷抵消，球壳外产生电场。但壳内电荷的位置并不影响外层电荷的分布。

\note{导体球壳的电势公式并不要求电荷是均匀分布的。}

\begin{example}{导体球与球壳}
    内外半径分别为$2R,3R$的导体球壳带有电荷$+Q$，球壳空腔内还有一个点电荷$+Q$，到球心距离为$R$。(1)~求球壳内外表面的电荷，它们的分布均匀吗？(2)~求球心处的电势。

    【补充练习】内外半径分别为$2R,3R$的导体球壳不带电，空腔内还有一个半径为$R$的同心金属球，带电量为$+Q$。(1)~求电荷、电场和电势的分布。(2)~若将球壳接地再断开，求电荷、电场和电势的分布；(3)~若将球壳与金属球用导线相连，求电荷、电场和电势的分布。
\end{example}

\section{电场中的电介质}

\subsection{电介质的电结构特点}

电介质就是通常所说的绝缘体，内部没有可自由移动的电荷。但在电场中，也会因电场的影响，发生电极化现象。

\begin{itemize}
    \item 极性分子在电场中，会发生取向的变化
    \item 非极性分子在电场中，会发生电荷分异
\end{itemize}

在外电场的作用下，电介质表面会束缚电荷（或称为极化电荷），这种现象称为电介质的极化。

\note{导体内部的电荷移动会完全抵消外电场，而电介质的极化只能部分抵消外电场。因此在电介质内部也有电场，但比外电场要弱一些。}

\subsection{电位移}

\subsubsection{电位移矢量}

有电介质存在时，电场分布可用电位移$\vec{D}$描述，单位是$\unit{C/m^2}$。

\begin{formula}
    $\vec{D}=\varepsilon \vec{E}$
\end{formula}

其中，$\varepsilon$表示电介质的介电常数（或电容率），真空的介电常数就是$\varepsilon_0$，因此真空的$\vec{D}=\varepsilon_0 \vec{E}$

介电常数还可以用相对介电常数$\varepsilon_r$描述，表示是真空的多少倍，则$\varepsilon=\varepsilon_0\varepsilon_r$

\note{这一关系仅对各向同性的电介质成立，大学物理课程仅考虑这种电介质。}

\subsubsection{电极化强度}

有电介质存在时，电场是电介质的束缚电荷和其他自由电荷共同产生的，因此电介质与真空的电位移之差，就描述了电介质的极化情况，称为电极化强度$\vec{P}$

\begin{formula}
    $\vec{P}=\vec{D}-\varepsilon_0\vec{E}=\varepsilon_0(\varepsilon_r-1)\vec{E}=\varepsilon_0\chi\vec{E}$
\end{formula}

其中，$\chi=\varepsilon_r-1$也称为电介质的电极化率。

\subsubsection{D的高斯定理}

有电介质存在时，高斯定理仍然成立，即$\displaystyle \varoiint_S {\vec{E}\cdot d\vec{S}}=\dfrac{1}{\varepsilon_0}\sum q_{\text{内}}$，但因为束缚电荷很难知道，所以并不实用。

无论有无电介质存在，通过任意封闭曲面的电位移通量，等于该封闭面包围的自由电荷的代数和，这就是D的高斯定理。

\begin{formula}
    $\displaystyle \varoiint_S {\vec{D}\cdot d\vec{S}}=\sum q_{\text{内自由}}$
\end{formula}

\note{若没有电介质存在，则$\vec{D}=\varepsilon_0 \vec{E}$，就划归为前面所学的高斯定理。}

\warn{注意D的高斯定理只包含内部的自由电荷，不包含束缚电荷。}

\begin{example}{D的高斯定理}
    (1)~若高斯面内不包围自由电荷，则高斯面的下列哪些量为零：(a)~任意点的电场强度；(b)~任意点的电位移矢量；(c)~E通量；(d)~D通量。(2)~若闭合曲面的D通量为零，说明面内电荷情况如何：(a)~没有束缚电荷；(b)~没有自由电荷；(c)~束缚电荷与自由电荷的代数和为零；(d)~自由电荷的代数和为零。
    
    【补充练习】判断以下说法是否正确：(1)~若选择的高斯面通过电介质，则D的高斯定理成立，E的高斯定理不成立；(2)~若高斯面内没有自由电荷，则D通量和E通量均为零；(3)~若高斯面上的电位移矢量处处为零，说明内部自由电荷的代数和为零；(4)~高斯面的D通量仅与内部的自由电荷有关，而E通量与内部和外部的自由电荷都有关。
\end{example}

\begin{example}{有电介质存在的场强}
    一个半径为$R$的薄金属球壳，均匀带电$+Q$。(1)~若球壳内部充满相对介电常数为$\varepsilon_r$的电介质，外部为真空，求电场强度和电势的分布；(2)~若球壳内部为真空，外部包了一层厚度为$R$、相对介电常数为$\varepsilon_r$的电介质，求电场强度和电势的分布。
    \tcblower
    \tcbline
    【补充练习】半径为$R_1,R_2$的无限长同轴金属圆筒，其间充满相对介电常数为$\varepsilon_r$的电介质，设内外圆筒单位长度所带电荷分别为$+\lambda,-\lambda$。求介质中到轴线距离为$r$的电场强度与电位移矢量的大小。
\end{example}

\section{电容}

\subsection{电容器及其电容}

\subsubsection{电容器}

两个彼此绝缘但靠近的导体形成电容器，其电容定义为

\begin{formula}
    $C=\dfrac{Q}{U}$
\end{formula}

\subsubsection{常见的电容器}

电容器的电容可用D的高斯定理求出，一些常见的电容器的电容如下表。

\begin{table}[htb]
    \centering
    \begin{tabular}{|c|c|c|}
        \hline
        电容器的形式 & 图示 & 电容\\
        \hline
        平行板电容器 &  & $C=\dfrac{\varepsilon S}{d}$\\
         \hline
        圆柱形电容器 &  & $C=\dfrac{2\pi \varepsilon L}{\ln{(R_2/R_1)}}$\\
         \hline
        球形电容器 &  & $C=\dfrac{4\pi \varepsilon R_1R_2}{R_2-R_1}$\\
         \hline
        孤立导体球 &  & $C=4\pi\varepsilon R$\\
         \hline
    \end{tabular}
\end{table}

\note{对于一个孤立导体，认为它与无限远处的另一导体组成电容器。}

\subsection{电容器的能量}

电容器的能量储存在电场之中，公式为

\begin{formula}
    $W=\dfrac{1}{2}QU=\dfrac{1}{2}CU^2=\dfrac{1}{2}\cdot\dfrac{Q^2}{C}$
\end{formula}

\begin{example}{电容器}
    内外半径分别为$R_1,R_2$的圆柱形电容器，间隙填满相对介电常数为$\varepsilon_r$的电介质，若电容内的电场强度最大不能超过$E_m$，求达到极限时的电压、电荷量、储存的电能。
    
    【补充练习】内外半径分别为$R_1,R_2$的同心球形电容器，间隙填满相对介电常数为$\varepsilon_r$的电介质。当内外球面带电量为$Q$时，求：(1)~内外球面的电势差；(2)~电场强度的分布；(3)~储存的电能。
\end{example}

\begin{example}{电容的动态分析}
    一平行板电容器，极板之间为真空，极板间距为$d$，面积为$S$，充电后切断电源。分别进行以下操作，判断电容、带电量、电场强度、电压、电场能量如何变化。(1)~将极板间距变为$2d$；(2)~将极板错开，使相对面积变为$S/2$；(3)~插入一块厚度为$d/2$的金属导体；(4)~填满相对介电常数为$\varepsilon_r$的电介质。
    \tcblower
    \tcbline
    【补充练习】一平行板电容器，极板之间为真空，极板间距为$d$，面积为$S$，始终与直流电源相连。分别进行以下操作，判断电容、带电量、电场强度、电压、电场能量如何变化。(1)~将极板间距变为$2d$；(2)~将极板错开，使相对面积变为$S/2$；(3)~插入一块厚度为$d/2$的金属导体；(4)~填满相对介电常数为$\varepsilon_r$的电介质。
    
\end{example}

\subsection{电容的串并联}

\begin{itemize}
    \item 几个电容并联，总电容等于各个电容之和
    \item 几个电容串联，总电容的倒数等于各个电容的倒数和
\end{itemize}

\begin{formula}
    $C_{\text{并}}=\sum{C_i},~\dfrac{1}{C_{\text{串}}}=\sum{\dfrac{1}{C_i}}$
\end{formula}

\note{电容串并联规律与电阻是相反的，与弹簧是相同的。}



\section{拓展内容}

\subsection{唯一性定理与镜像法}



\subsection{静电场的边界条件}

\subsection{电场的能量}

\subsubsection{电荷系的能量}

对于多个电荷组成的系统，把各个电荷从无穷远聚集到现在位置，克服电场力所做的功，称为电荷系的静电能。静电能也可以定义为，将各电荷从现有位置分散到无穷远，静电力所做的功。

\begin{formula}
    离散：$W=\dfrac{1}{2}\sum{q_i\varphi_i}$~~~连续：$W=\dfrac{1}{2}\displaystyle\int_q{\varphi dq}$
\end{formula}

其中，$\varphi_i$表示除$q_i$外的其他电荷在$q_i$处产生的电势。

\subsubsection{电场的能量}

电场在单位体积内的能量称为能量密度$w$，电场的总能量$W$就是全空间上能量密度的体积分。

\begin{formula}
    $w=\dfrac{1}{2}\vec{D}\cdot\vec{E}=\dfrac{1}{2}\varepsilon E^2,~~W=\displaystyle\iiint_V wdV$
\end{formula}

\begin{example}{电场的能量}
    有一橡胶做的气球，可视为半径为$R$的球面，电荷$+Q$在上面均匀分布。(1)~求电场的能量；(2)~若电荷量变为原来的2倍，电场能量如何变化？(3)~若电荷量不变，但将其吹大，半径从$R$变为$2R$，则电场能量如何变化？
    
    【补充练习】地球表面上空晴天时的电场强度约为$\SI{100}{V/m}$。(1)~求电场的能量密度；(2)~若整个地球表面以上$\SI{10}{km}$范围内电场强度都是这一数值，取地球半径为$\SI{6400}{km}$，求这一范围内的电场能。
\end{example}

\subsection{电容的串并联分析}

\begin{example}{电容的串并联}
    如图所示，空气电容器$C_1=\SI{10}{\mu F},C_2=C_3=\SI{5}{\mu F}$。(1)~求总电容；(2)~若连接$\SI{100}{V}$直流电源充电，求各电容的电荷量和电压；(3)~若保持连接直流电源，将一电介质插入$C_1$中，各电容器的电压如何变化？(4)~若将电源断开，再将一电介质插入$C_1$中，各电容器的电压如何变化？
    
    【补充练习】一平行板电容器，极板之间为真空，极板间距为$d$，面积为$S$，充电后切断电源。在其中插入一块电介质，两极板的电压变为原来的一半，在以下两种情况下求相对介电常数，并求出自由电荷与极化电荷的分布。(1)~电介质的面积也为$S$，厚度为$d/3$；(2)~电介质的厚度为$d$，但面积为$S/3$.
\end{example}